The angular diameter of the Sun during perihelion is
\[
d_1=2\times\frac{695{,}500}{147{,}100{,}00-6{,}378}
=2\times0.004731=0.009463\ \mathrm{rad},
\]
% sic: the printed denominator is 147,100,00, missing one zero compared with the question's 147,100,000 km.
while the angular diameter of the Moon during apogee is
\[
d_2=2\times\frac{1{,}737}{405{,}000-6{,}378}
=2\times0.004358=0.008715\ \mathrm{rad}.
\]
\hfill\mbox{(6 points)}

The fraction of the area covered is the ratio of the squares of these values,
\[
r=\left(\frac{d_2}{d_1}\right)^2=0.8482=84.82\%.
\]
\hfill\mbox{(4 points)}
